\chapter{Pages web et formulaires (rappels)}

\begin{synthese}[Au programme]
Avant de programmer côté serveur en PHP, il faut maîtriser le socle qui s'exécute
dans le navigateur de l'élève : le \textbf{HTML} (structure de la page), les
\textbf{formulaires} (saisie de l'utilisateur) et le \textbf{JavaScript} (vérification
de la saisie). C'est la base de tout site web dynamique.
\end{synthese}

\section{Structure d'une page HTML}

\begin{definition}
\textbf{HTML} (\emph{HyperText Markup Language}) est le langage de \textbf{balisage}
qui décrit la structure d'une page web. Une balise s'écrit entre chevrons : une
balise \textbf{ouvrante} \code{<p>} et une balise \textbf{fermante} \code{</p>}
entourent un contenu. Le navigateur lit le HTML et l'\textbf{affiche}.
\end{definition}

Toute page commence par une déclaration \code{<!DOCTYPE html>} qui indique qu'il
s'agit d'une page HTML moderne. Le document se divise en deux parties :
le \code{<head>} (informations sur la page : titre, encodage) et le \code{<body>}
(le contenu visible).

\begin{codehtml}
<!DOCTYPE html>
<html lang="fr">
<head>
  <meta charset="utf-8">           <!-- accents : é è à ç -->
  <title>Lycée Bilingue de Bafoussam</title>
</head>
<body>
  <h1>Bienvenue au lycée</h1>       <!-- titre principal -->
  <h2>Classe de Terminale TI</h2>   <!-- sous-titre -->
  <p>Cette page presente nos eleves.</p>  <!-- paragraphe -->
  <a href="cours.html">Voir les cours</a> <!-- lien -->
</body>
</html>
\end{codehtml}

\fig{Squelette minimal d'une page HTML : doctype, en-tête et corps.}

\begin{propriete}[Balises courantes]
\code{<h1>} à \code{<h6>} : titres (du plus grand au plus petit) ; \code{<p>} :
paragraphe ; \code{<br>} : saut de ligne ; \code{<a href="...">} : lien ;
\code{<img src="..." alt="...">} : image ; \code{<ul>}/\code{<ol>} avec \code{<li>} :
listes ; \code{<div>} et \code{<span>} : conteneurs neutres.
\end{propriete}

\begin{attention}
Certaines balises sont \textbf{vides} (sans contenu) et ne se ferment pas :
\code{<br>}, \code{<img>}, \code{<input>}, \code{<meta>}. Les autres doivent
\textbf{toujours} être refermées : à chaque \code{<p>} correspond un \code{</p>}.
\end{attention}

\section{Les formulaires HTML}

Un formulaire permet à l'utilisateur (un élève, un parent\dots) de \textbf{saisir}
des informations puis de les \textbf{envoyer} au serveur. C'est par lui que les
données entrent dans une application web.

\subsection{La balise \texttt{<form>}}

\begin{syntaxe}
\code{<form action="traitement.php" method="post"> \dots champs \dots </form>}
\end{syntaxe}

\begin{itemize}
  \item \code{action} : l'adresse du fichier (souvent un script PHP) qui va
        \textbf{traiter} les données envoyées ;
  \item \code{method} : la façon d'envoyer les données. \code{post} (les données
        sont cachées, recommandé pour les formulaires) ou \code{get} (les données
        apparaissent dans l'URL).
\end{itemize}

\begin{attention}
Chaque champ destiné à être envoyé \textbf{doit} posséder un attribut \code{name}.
C'est ce nom que le serveur (PHP) utilisera pour récupérer la valeur. Un champ
\textbf{sans} \code{name} n'est pas transmis.
\end{attention}

\subsection{Les champs de saisie : \texttt{<input>}}

La balise \code{<input>} change de comportement selon son attribut \code{type}.

\begin{center}\small
\begin{tabular}{@{}ll@{}}
\toprule
\textbf{type} & \textbf{Rôle}\\
\midrule
\code{text} & saisie d'une ligne de texte (nom, ville\dots)\\
\code{password} & saisie masquée (mot de passe : \texttt{••••})\\
\code{number} & nombre uniquement (âge, note)\\
\code{email} & adresse e-mail (format vérifié par le navigateur)\\
\code{radio} & bouton radio : un seul choix dans un groupe\\
\code{checkbox} & case à cocher : choix multiples\\
\code{submit} & bouton qui envoie le formulaire\\
\bottomrule
\end{tabular}
\end{center}

\begin{codehtml}
<form action="inscription.php" method="post">

  <!-- Champ texte avec une étiquette cliquable -->
  <label for="nom">Nom de l'élève :</label>
  <input type="text" id="nom" name="nom">
  <br>

  <!-- Mot de passe (saisie masquée) -->
  <label for="mdp">Mot de passe :</label>
  <input type="password" id="mdp" name="mdp">
  <br>

  <!-- Nombre : l'âge -->
  <label for="age">Âge :</label>
  <input type="number" id="age" name="age">
  <br>

  <!-- Adresse e-mail -->
  <label for="mail">E-mail :</label>
  <input type="email" id="mail" name="mail">
  <br>

  <!-- Boutons radio : un SEUL choix (même name) -->
  Sexe :
  <input type="radio" name="sexe" value="M"> Masculin
  <input type="radio" name="sexe" value="F"> Féminin
  <br>

  <!-- Case à cocher -->
  <input type="checkbox" name="interne" value="oui"> Élève interne
  <br>

  <!-- Bouton d'envoi -->
  <input type="submit" value="S'inscrire">
</form>
\end{codehtml}

\fig{Un formulaire d'inscription complet utilisant plusieurs types de champs.}

\begin{remarque}
La balise \code{<label for="nom">} est reliée au champ dont l'\code{id} vaut
\code{nom}. Cliquer sur l'étiquette place alors le curseur dans le champ : c'est
plus pratique, surtout sur téléphone.
\end{remarque}

\subsection{Liste déroulante : \texttt{<select>}}

Pour proposer un choix dans une liste fixe (la série, la classe\dots), on utilise
\code{<select>} qui contient des \code{<option>}.

\begin{codehtml}
<label for="serie">Série :</label>
<select id="serie" name="serie">
  <option value="TI">Terminale TI</option>
  <option value="C">Terminale C</option>
  <option value="D">Terminale D</option>
</select>
\end{codehtml}

\fig{La valeur envoyée est celle de l'attribut \code{value} de l'option choisie.}

\subsection{Zone de texte : \texttt{<textarea>}}

Pour un texte long (une remarque, une adresse complète), on emploie
\code{<textarea>}, qui s'écrit sur plusieurs lignes.

\begin{codehtml}
<label for="obs">Observations :</label><br>
<textarea id="obs" name="obs" rows="4" cols="40">
Écrivez ici...
</textarea>
\end{codehtml}

\section{Les tableaux HTML}

\begin{definition}
Un \textbf{tableau} HTML organise des données en lignes et colonnes. Il se construit
avec \code{<table>} (le tableau), \code{<tr>} (\emph{table row} : une ligne),
\code{<th>} (\emph{table header} : cellule d'en-tête) et \code{<td>}
(\emph{table data} : cellule ordinaire).
\end{definition}

\begin{codehtml}
<table border="1">
  <!-- Première ligne : les en-têtes -->
  <tr>
    <th>Matière</th>
    <th>Note /20</th>
  </tr>
  <!-- Lignes de données -->
  <tr>
    <td>Informatique</td>
    <td>16</td>
  </tr>
  <tr>
    <td>Mathématiques</td>
    <td>13</td>
  </tr>
</table>
\end{codehtml}

\fig{Un tableau de notes : une ligne d'en-tête (\code{<th>}) puis les données (\code{<td>}).}

\begin{remarque}
L'attribut \code{border="1"} dessine une bordure simple. Dans un vrai site, on
préfère colorer et espacer les cellules avec du \textbf{CSS}, mais \code{border}
suffit pour visualiser la structure pendant l'apprentissage.
\end{remarque}

\section{JavaScript : les bases}

\begin{definition}
\textbf{JavaScript} (abrégé \textbf{JS}) est un langage de programmation qui
s'exécute \textbf{dans le navigateur}, donc \textbf{côté client}. Il sert à rendre
la page \emph{interactive} : réagir à un clic, vérifier un formulaire, modifier le
contenu affiché.
\end{definition}

On insère du JavaScript dans la balise \code{<script>}, généralement juste avant
\code{</body>}.

\begin{syntaxe}
\code{<script>}\quad \textit{instructions JavaScript\dots}\quad \code{</script>}
\end{syntaxe}

\subsection{Variables et types}

On déclare une variable avec \code{let} (modifiable), \code{const} (constante,
non modifiable) ou l'ancien mot-clé \code{var}. JavaScript devine le type tout seul.

\begin{codejs}
<script>
  let nom = "Awa";         // chaîne de caractères
  let age = 17;            // nombre
  const PI = 3.14;         // constante (ne change plus)
  let estAdmis = true;     // booléen (true ou false)
  // Afficher dans la console du navigateur
  console.log(nom, age, estAdmis);
</script>
\end{codejs}

\begin{attention}
En JavaScript, chaque instruction se termine par un point-virgule \code{;} et les
mots-clés s'écrivent en \textbf{minuscules}. \code{Let} ou \code{LET} provoquent
une erreur : il faut écrire \code{let}.
\end{attention}

\subsection{Opérateurs}

\begin{center}\small
\begin{tabular}{@{}lll@{}}
\toprule
\textbf{Catégorie} & \textbf{Opérateurs} & \textbf{Exemple}\\
\midrule
Arithmétiques & \code{+ - * / \%} & \code{age + 1}\\
Comparaison & \code{== != < > <= >=} & \code{note >= 10}\\
Égalité stricte & \code{=== !==} & \code{type === "number"}\\
Logiques & \code{\&\& || !} & \code{x>0 \&\& x<20}\\
\bottomrule
\end{tabular}
\end{center}

\begin{remarque}
Le \code{+} sert aussi à \textbf{coller} deux chaînes :
\code{"Bon" + "jour"} donne \code{"Bonjour"}. C'est la \emph{concaténation}.
\end{remarque}

\subsection{Fonctions, condition et boucle}

Une \textbf{fonction} regroupe des instructions sous un nom, pour les réutiliser.
Elle peut \textbf{retourner} un résultat avec \code{return}.

\begin{codejs}
<script>
  // Fonction qui calcule une moyenne de deux notes
  function moyenne(n1, n2) {
    return (n1 + n2) / 2;
  }

  let m = moyenne(12, 16);   // m vaut 14

  // Condition : if ... else
  if (m >= 10) {
    console.log("Admis avec " + m);
  } else {
    console.log("Recalé");
  }

  // Boucle for : compter de 1 à 5
  for (let i = 1; i <= 5; i = i + 1) {
    console.log("Élève numéro " + i);
  }
</script>
\end{codejs}

\fig{Fonction, condition \code{if/else} et boucle \code{for} : les briques de base.}

\section{JavaScript et la page}

JavaScript devient vraiment utile quand il \textbf{réagit} aux actions de
l'utilisateur et \textbf{lit} le contenu de la page.

\subsection{Les événements}

Un \textbf{événement} est une action de l'utilisateur (clic, envoi de formulaire).
On y associe du code grâce à des attributs HTML :

\begin{itemize}
  \item \code{onclick} : se déclenche au clic sur un élément (souvent un bouton) ;
  \item \code{onsubmit} : se déclenche quand on envoie un formulaire.
\end{itemize}

\subsection{Lire un champ et afficher un message}

Pour récupérer ce que l'élève a tapé, on retrouve le champ par son \code{id}, puis
on lit sa propriété \code{.value}. La fonction \code{alert(\dots)} affiche, elle,
une petite fenêtre d'information.

\begin{syntaxe}
\code{document.getElementById("id").value} renvoie le texte saisi dans le champ
dont l'attribut \code{id} vaut \code{"id"}.
\end{syntaxe}

\begin{codehtml}
<input type="text" id="ville" name="ville">
<button onclick="direBonjour()">Saluer</button>

<script>
  function direBonjour() {
    // On lit la valeur tapée dans le champ "ville"
    let v = document.getElementById("ville").value;
    alert("Bonjour, élève de " + v + " !");
  }
</script>
\end{codehtml}

\fig{Au clic, la fonction lit le champ \code{ville} et affiche un message avec \code{alert}.}

\begin{methode}
\textbf{Tester du JavaScript.} (1)~Enregistre le fichier avec l'extension
\code{.html} ; (2)~ouvre-le directement dans un navigateur (Chrome, Firefox) ;
(3)~ouvre la \textbf{console} (touche \code{F12}) pour voir les messages
\code{console.log} et les éventuelles erreurs.
\end{methode}

\section{Contrôle de saisie d'un formulaire}

\begin{definition}
Le \textbf{contrôle de saisie} (ou \emph{validation}) consiste à vérifier, avec
JavaScript, que les informations tapées par l'utilisateur sont \textbf{correctes}
\textbf{avant} de les envoyer au serveur. Cela évite d'enregistrer des données
fausses ou vides.
\end{definition}

\subsection{Le principe : bloquer l'envoi}

On relie une fonction de vérification à l'événement \code{onsubmit} du formulaire.
\textbf{Règle d'or} : si la fonction renvoie \code{false}, le formulaire
\textbf{n'est pas envoyé} ; si elle renvoie \code{true}, l'envoi continue.

\begin{syntaxe}
\code{<form onsubmit="return verifier()" \dots>} : l'envoi est annulé quand
\code{verifier()} renvoie \code{false}.
\end{syntaxe}

\subsection{Les vérifications utiles}

\begin{itemize}
  \item \textbf{Champ vide} : comparer la valeur à la chaîne vide
        (\code{valeur === ""}) ;
  \item \textbf{Valeur numérique} : convertir avec \code{Number(valeur)} et tester
        avec \code{isNaN(...)} (\emph{is Not a Number}, « n'est pas un nombre ») ;
  \item \textbf{Longueur} : la propriété \code{.length} donne le nombre de
        caractères (\code{mdp.length < 6}) ;
  \item \textbf{Format simple} : chercher un caractère, par exemple le \code{@}
        d'un e-mail avec \code{valeur.indexOf("@")} (vaut \code{-1} s'il est absent).
\end{itemize}

\subsection{Exemple complet}

Voici un formulaire d'inscription complet avec sa fonction de validation. La
fonction vérifie successivement chaque champ et renvoie \code{false} dès qu'un
problème est détecté, ce qui \textbf{bloque} l'envoi.

\begin{codehtml}
<!DOCTYPE html>
<html lang="fr">
<head><meta charset="utf-8"><title>Inscription</title></head>
<body>

<h2>Inscription — Lycée de Garoua</h2>

<!-- L'envoi n'a lieu que si verifier() renvoie true -->
<form action="inscription.php" method="post"
      onsubmit="return verifier()">

  <label for="nom">Nom :</label>
  <input type="text" id="nom" name="nom"><br>

  <label for="age">Âge :</label>
  <input type="text" id="age" name="age"><br>

  <label for="mdp">Mot de passe :</label>
  <input type="password" id="mdp" name="mdp"><br>

  <input type="submit" value="S'inscrire">
</form>

<script>
  function verifier() {
    // 1) Lire les valeurs tapées par l'élève
    let nom = document.getElementById("nom").value;
    let age = document.getElementById("age").value;
    let mdp = document.getElementById("mdp").value;

    // 2) Le nom ne doit pas être vide
    if (nom === "") {
      alert("Le nom est obligatoire.");
      return false;          // on bloque l'envoi
    }

    // 3) L'âge doit être un nombre strictement positif
    if (isNaN(age) || age === "") {
      alert("L'âge doit être un nombre.");
      return false;
    }
    if (Number(age) <= 0) {
      alert("L'âge doit être positif.");
      return false;
    }

    // 4) Le mot de passe doit faire au moins 6 caractères
    if (mdp.length < 6) {
      alert("Mot de passe : 6 caractères minimum.");
      return false;
    }

    // 5) Tout est correct : on autorise l'envoi
    return true;
  }
</script>

</body>
</html>
\end{codehtml}

\fig{Formulaire + fonction \code{verifier()} : chaque \code{return false} empêche l'envoi.}

\begin{attention}
La validation JavaScript améliore le confort de l'élève, mais elle s'exécute dans
\textbf{son} navigateur : un utilisateur malintentionné peut la contourner. Les
données devront \textbf{toujours} être revérifiées côté serveur (en PHP, au
chapitre sur les formulaires).
\end{attention}

\section{Exercices}

\rubrique{Application}

\exo{Page de présentation}\diff{1}
Écris une page HTML complète (doctype, head, body) qui affiche le titre
\emph{« Lycée Technique de Douala »}, un sous-titre \emph{« Terminale TI »} et un
paragraphe de présentation de ton choix.

\exo{Formulaire d'inscription d'un élève}\diff{1}
Crée un formulaire (méthode \code{post}, action \code{inscrire.php}) avec : un champ
texte \emph{Nom}, un champ texte \emph{Prénom}, un champ \code{number} \emph{Âge},
une liste déroulante \emph{Série} (TI, C, D) et un bouton d'envoi. Chaque champ doit
avoir un \code{name} et un \code{label}.

\exo{Tableau des notes}\diff{1}
Construis un tableau HTML à deux colonnes (\emph{Matière}, \emph{Note /20}) et trois
lignes de données : Informatique 15, Français 11, Physique 13. Ajoute une ligne
d'en-tête avec \code{<th>}.

\exo{Choix multiples}\diff{2}
Réalise un formulaire demandant le \emph{sexe} (deux boutons radio M/F) et les
\emph{options choisies} (cases à cocher : Sport, Musique, Club informatique).
Termine par un bouton d'envoi.

\exo{Saluer l'élève}\diff{2}
Crée un champ texte (\code{id="prenom"}) et un bouton. Au clic sur le bouton, une
fonction JavaScript lit le prénom et affiche, avec \code{alert}, le message
\emph{« Bienvenue, \dots ! »}.

\rubrique{Entraînement}

\exo{Validation de l'âge}\diff{2}
À partir d'un formulaire contenant un champ \emph{Âge}, écris une fonction
JavaScript reliée à \code{onsubmit} qui bloque l'envoi si l'âge n'est pas un nombre
ou s'il n'est pas strictement positif.

\exo{Deux mots de passe identiques}\diff{2}
Crée un formulaire avec deux champs \code{password} (\emph{Mot de passe} et
\emph{Confirmation}). Écris une fonction qui bloque l'envoi si les deux valeurs sont
différentes, ou si le mot de passe fait moins de 6 caractères.

\exo{Calculer une moyenne}\diff{2}
Trois champs \code{number} contiennent trois notes. Au clic sur un bouton, une
fonction JavaScript calcule la moyenne et l'affiche avec \code{alert}, en précisant
\emph{« Admis »} si la moyenne est $\geq 10$, sinon \emph{« Recalé »}.

\exo{E-mail valide (format simple)}\diff{3}
Dans un formulaire d'inscription, ajoute un champ \emph{e-mail}. Écris une fonction
de validation qui bloque l'envoi si le champ est vide ou s'il ne contient pas le
caractère \code{@} (utilise \code{indexOf}).

\exo{Inscription complète vérifiée}\diff{3}
Reprends le formulaire de l'exercice 2 et écris une seule fonction \code{verifier()}
qui contrôle : nom non vide, prénom non vide, âge numérique et positif. L'envoi ne
doit avoir lieu que si tout est correct.

\section{Corrigés}

\corrige{de l'exercice 1}
\begin{codehtml}
<!DOCTYPE html>
<html lang="fr">
<head>
  <meta charset="utf-8">
  <title>Lycée Technique de Douala</title>
</head>
<body>
  <h1>Lycée Technique de Douala</h1>
  <h2>Terminale TI</h2>
  <p>Notre classe forme les futurs techniciens en informatique
     de la ville de Douala.</p>
</body>
</html>
\end{codehtml}

\corrige{de l'exercice 2}
\begin{codehtml}
<form action="inscrire.php" method="post">
  <label for="nom">Nom :</label>
  <input type="text" id="nom" name="nom"><br>

  <label for="prenom">Prénom :</label>
  <input type="text" id="prenom" name="prenom"><br>

  <label for="age">Âge :</label>
  <input type="number" id="age" name="age"><br>

  <label for="serie">Série :</label>
  <select id="serie" name="serie">
    <option value="TI">Terminale TI</option>
    <option value="C">Terminale C</option>
    <option value="D">Terminale D</option>
  </select><br>

  <input type="submit" value="S'inscrire">
</form>
\end{codehtml}

\corrige{de l'exercice 3}
\begin{codehtml}
<table border="1">
  <tr>
    <th>Matière</th>
    <th>Note /20</th>
  </tr>
  <tr><td>Informatique</td><td>15</td></tr>
  <tr><td>Français</td><td>11</td></tr>
  <tr><td>Physique</td><td>13</td></tr>
</table>
\end{codehtml}

\corrige{de l'exercice 4}
\begin{codehtml}
<form action="choix.php" method="post">
  Sexe :
  <input type="radio" name="sexe" value="M"> Masculin
  <input type="radio" name="sexe" value="F"> Féminin
  <br>
  Options :
  <input type="checkbox" name="sport" value="oui"> Sport
  <input type="checkbox" name="musique" value="oui"> Musique
  <input type="checkbox" name="club" value="oui"> Club informatique
  <br>
  <input type="submit" value="Valider">
</form>
\end{codehtml}

\corrige{de l'exercice 5}
\begin{codehtml}
<input type="text" id="prenom" name="prenom">
<button onclick="saluer()">Saluer</button>

<script>
  function saluer() {
    // Lecture du prénom tapé
    let p = document.getElementById("prenom").value;
    alert("Bienvenue, " + p + " !");
  }
</script>
\end{codehtml}

\corrige{de l'exercice 6}
\begin{codehtml}
<form action="age.php" method="post" onsubmit="return verifierAge()">
  <label for="age">Âge :</label>
  <input type="text" id="age" name="age">
  <input type="submit" value="Envoyer">
</form>

<script>
  function verifierAge() {
    let age = document.getElementById("age").value;
    // Doit être un nombre
    if (isNaN(age) || age === "") {
      alert("L'âge doit être un nombre.");
      return false;
    }
    // Doit être strictement positif
    if (Number(age) <= 0) {
      alert("L'âge doit être positif.");
      return false;
    }
    return true;   // tout est correct : envoi autorisé
  }
</script>
\end{codehtml}

\corrige{de l'exercice 7}
\begin{codehtml}
<form action="compte.php" method="post" onsubmit="return verifierMdp()">
  <label for="mdp">Mot de passe :</label>
  <input type="password" id="mdp" name="mdp"><br>

  <label for="conf">Confirmation :</label>
  <input type="password" id="conf" name="conf"><br>

  <input type="submit" value="Créer le compte">
</form>

<script>
  function verifierMdp() {
    let mdp = document.getElementById("mdp").value;
    let conf = document.getElementById("conf").value;

    // Longueur minimale
    if (mdp.length < 6) {
      alert("Le mot de passe doit faire au moins 6 caractères.");
      return false;
    }
    // Les deux saisies doivent être identiques
    if (mdp !== conf) {
      alert("Les deux mots de passe sont différents.");
      return false;
    }
    return true;
  }
</script>
\end{codehtml}

\corrige{de l'exercice 8}
\begin{codehtml}
N1 : <input type="number" id="n1"><br>
N2 : <input type="number" id="n2"><br>
N3 : <input type="number" id="n3"><br>
<button onclick="calculer()">Calculer la moyenne</button>

<script>
  function calculer() {
    // On convertit chaque saisie en nombre
    let n1 = Number(document.getElementById("n1").value);
    let n2 = Number(document.getElementById("n2").value);
    let n3 = Number(document.getElementById("n3").value);

    let moy = (n1 + n2 + n3) / 3;

    if (moy >= 10) {
      alert("Moyenne : " + moy + " — Admis");
    } else {
      alert("Moyenne : " + moy + " — Recalé");
    }
  }
</script>
\end{codehtml}

\corrige{de l'exercice 9}
\begin{codehtml}
<form action="inscrire.php" method="post" onsubmit="return verifierMail()">
  <label for="mail">E-mail :</label>
  <input type="text" id="mail" name="mail">
  <input type="submit" value="Envoyer">
</form>

<script>
  function verifierMail() {
    let mail = document.getElementById("mail").value;

    // Champ vide ?
    if (mail === "") {
      alert("L'e-mail est obligatoire.");
      return false;
    }
    // Le caractère @ est-il présent ? indexOf vaut -1 sinon
    if (mail.indexOf("@") === -1) {
      alert("E-mail invalide : le caractère @ est manquant.");
      return false;
    }
    return true;
  }
</script>
\end{codehtml}

\corrige{de l'exercice 10}
\begin{codehtml}
<form action="inscrire.php" method="post" onsubmit="return verifier()">
  <label for="nom">Nom :</label>
  <input type="text" id="nom" name="nom"><br>

  <label for="prenom">Prénom :</label>
  <input type="text" id="prenom" name="prenom"><br>

  <label for="age">Âge :</label>
  <input type="text" id="age" name="age"><br>

  <input type="submit" value="S'inscrire">
</form>

<script>
  function verifier() {
    let nom = document.getElementById("nom").value;
    let prenom = document.getElementById("prenom").value;
    let age = document.getElementById("age").value;

    // Nom et prénom obligatoires
    if (nom === "") {
      alert("Le nom est obligatoire.");
      return false;
    }
    if (prenom === "") {
      alert("Le prénom est obligatoire.");
      return false;
    }
    // Âge : un nombre strictement positif
    if (isNaN(age) || age === "") {
      alert("L'âge doit être un nombre.");
      return false;
    }
    if (Number(age) <= 0) {
      alert("L'âge doit être positif.");
      return false;
    }
    return true;   // formulaire correct : envoi autorisé
  }
</script>
\end{codehtml}

\begin{aretenir}
Le \textbf{HTML} structure la page, le \textbf{formulaire} (\code{<form>}, champs
\code{<input>}, attribut \code{name}) recueille la saisie, et le \textbf{JavaScript}
la \textbf{contrôle} dans le navigateur. La règle clé du contrôle de saisie :
\code{onsubmit="return verifier()"} et un \code{return false} pour
\textbf{bloquer} l'envoi tant qu'un champ est invalide. Mais ce contrôle reste
côté client : on revérifiera toujours côté serveur.
\end{aretenir}
